\chapter{Conclusion}

This Software Requirements Specification has defined QCanvas as a unified quantum simulation platform centered on OpenQASM 3.0. It established clear project boundaries, detailed functional and non-functional requirements, described the hybrid CPU–QPU execution model, and documented the system’s architecture, interfaces, and data design. Together, these specifications form the contract for development and verification, ensuring the system remains simulation-first, standards-aligned, and focused on educational and research use cases.

\section{Future Work}

While Iteration 1 and 2 have delivered a robust core platform, several key features are planned for the next development phase:

\begin{itemize}
    \item \textbf{Hybrid CPU-QPU Integration:} Validation and merging of the QCanvas orchestration layer with the QSim execution engine (Iteration 4).
    \item \textbf{Real-time WebSocket Communication:} Implementation of full duplex communication for live progress tracking of long-running simulations and batch conversions.
    \item \textbf{Advanced Quantum Simulation:} Integration of noise models and error mitigation techniques into the simulation backend.
    \item \textbf{User Authentication and Persistence:} Development of a secure user management system with database integration for saving projects and preferences.
    \item \textbf{Performance Optimization:} Benchmarking and optimizing the conversion and simulation pipelines for larger circuits.
\end{itemize}

This sets the stage for the final release of QCanvas as a comprehensive educational and research tool.
